\documentclass[12pt,a4paper,openany]{book}

% ------------------------------------------------------------
% Encodage, langue et typographie
% ------------------------------------------------------------
\usepackage[utf8]{inputenc}
\usepackage[T1]{fontenc}
\usepackage[french]{babel}
\usepackage{newpxtext}
\usepackage[scaled=0.94]{sourcesanspro}
\usepackage{microtype}
\usepackage{setspace}
\setstretch{1.18}
\usepackage{geometry}
\geometry{a4paper,inner=31mm,outer=25mm,top=27mm,bottom=29mm,bindingoffset=4mm}
\setlength{\parindent}{0pt}
\setlength{\parskip}{0.45em}
\raggedbottom
\clubpenalty=10000
\widowpenalty=10000
\displaywidowpenalty=10000
\emergencystretch=2em
\setcounter{tocdepth}{1}
\setcounter{secnumdepth}{1}

% ------------------------------------------------------------
% Mathématiques et tableaux
% ------------------------------------------------------------
\usepackage{amsmath,amssymb,amsthm,mathtools}
\usepackage{newpxmath}
\usepackage{booktabs,tabularx,array,multicol,graphicx}
\usepackage{caption}
\captionsetup{font={small,sf},labelfont={bf,color=navy},textfont={color=slate},skip=6pt}
\usepackage{enumitem}
\setlist[itemize]{leftmargin=1.35em,itemsep=0.2em,topsep=0.25em}
\setlist[enumerate]{leftmargin=1.55em,itemsep=0.25em,topsep=0.25em}
\newcolumntype{Y}{>{\raggedright\arraybackslash}X}

% ------------------------------------------------------------
% Graphiques : palette moderne et discrète
% ------------------------------------------------------------
\usepackage{xcolor}
\definecolor{ink}{HTML}{202B36}
\definecolor{navy}{HTML}{203A5F}
\definecolor{teal}{HTML}{2E6F78}
\definecolor{sage}{HTML}{55777A}
\definecolor{gold}{HTML}{A98245}
\definecolor{terracotta}{HTML}{A45C50}
\definecolor{slate}{HTML}{5E6B76}
\definecolor{mistblue}{HTML}{F5F7F9}
\definecolor{mistteal}{HTML}{F3F7F6}
\definecolor{mistsage}{HTML}{F4F7F6}
\definecolor{mistgold}{HTML}{FAF8F2}
\definecolor{mistred}{HTML}{FBF5F3}
\definecolor{softgray}{HTML}{F7F7F6}
% Compatibilité avec les figures déjà écrites.
\colorlet{uqamblue}{navy}
\colorlet{deepblue}{ink}
\colorlet{softblue}{mistblue}
\colorlet{softgreen}{mistsage}
\colorlet{softyellow}{mistgold}
\colorlet{softred}{mistred}
\usepackage{tikz}
\usetikzlibrary{arrows.meta,calc,positioning,patterns,angles,quotes,decorations.pathreplacing,3d}
\usepackage{pgfplots}
\pgfplotsset{compat=1.18}

% ------------------------------------------------------------
% Repères pédagogiques : surfaces légères, accents cohérents
% ------------------------------------------------------------
\usepackage[most]{tcolorbox}
\tcbset{
  enhanced jigsaw,
  breakable,
  frame hidden,
  boxrule=0pt,
  arc=0.8pt,
  outer arc=0.8pt,
  left=10pt,
  right=8pt,
  top=6pt,
  bottom=6pt,
  before skip=0.8em,
  after skip=0.8em,
  fonttitle=\sffamily\bfseries\small,
  separator sign none,
  before title={\vphantom{É}},
  after title={\par\smallskip}
}

\newtcolorbox{questionbox}[1][]{
  title={Question directrice},
  colback=mistteal,
  coltitle=teal,
  borderline west={2.2pt}{0pt}{teal},
  fontupper=\itshape,
  #1
}
\newtcolorbox{intuitionbox}[1][]{
  title={Idée intuitive},
  colback=mistsage,
  coltitle=sage,
  borderline west={1.8pt}{0pt}{sage},
  #1
}
\newtcolorbox{visualbox}[1][]{
  title={Lecture visuelle},
  colback=white,
  coltitle=teal,
  borderline west={2.1pt}{0pt}{teal},
  borderline north={0.35pt}{0pt}{teal!30},
  borderline south={0.35pt}{0pt}{teal!30},
  #1
}
\newtcolorbox{connectionbox}[1][]{
  title={Lien entre les représentations},
  colback=mistgold,
  coltitle=gold!75!black,
  borderline west={2pt}{0pt}{gold},
  #1
}
\newtcolorbox{methodbox}[1][]{
  title={Méthode},
  colback=mistblue,
  coltitle=navy,
  borderline west={2.2pt}{0pt}{navy},
  #1
}
\newtcolorbox{warningbox}[1][]{
  title={Point de vigilance},
  colback=mistred,
  coltitle=terracotta,
  borderline west={2.4pt}{0pt}{terracotta},
  #1
}
\newtcolorbox{summarybox}[1][]{
  title={Synthèse},
  colback=mistgold,
  coltitle=gold!75!black,
  borderline west={2.2pt}{0pt}{gold},
  #1
}
\newtcolorbox{examplebox}[1][]{
  title={Exemple},
  colback=softgray,
  coltitle=slate,
  borderline west={1.5pt}{0pt}{slate},
  #1
}
\newtcolorbox{counterbox}[1][]{
  title={Contre-exemple},
  colback=mistred,
  coltitle=terracotta,
  borderline west={2pt}{0pt}{terracotta},
  #1
}

\newtcbtheorem[number within=chapter]{definition}{Définition}{
  colback=mistteal,
  coltitle=teal,
  borderline west={2pt}{0pt}{teal},
  fonttitle=\sffamily\bfseries
}{def}
\newtcbtheorem[use counter from=definition]{theorem}{Théorème}{
  colback=mistblue,
  coltitle=navy,
  borderline west={2.4pt}{0pt}{navy},
  fonttitle=\sffamily\bfseries
}{thm}
\newtcbtheorem[use counter from=definition]{proposition}{Proposition}{
  colback=mistgold,
  coltitle=gold!75!black,
  borderline west={2pt}{0pt}{gold},
  fonttitle=\sffamily\bfseries
}{prop}

% ------------------------------------------------------------
% Titres, en-têtes et hyperliens
% ------------------------------------------------------------
\usepackage{titlesec}
\titleformat{\part}[display]
  {\centering\sffamily\bfseries\color{navy}}
  {\color{gold}\fontsize{44}{46}\selectfont\thepart}
  {0.6ex}
  {\Huge}
  [\vspace{1.1ex}{\color{teal}\titlerule}]
\titlespacing*{\part}{0pt}{1.5cm}{2.2cm}
\titleformat{\chapter}[display]
  {\sffamily\bfseries\color{navy}}
  {\raggedleft\color{teal}\large Chapitre \thechapter}
  {0.45ex}
  {\color{navy}\huge\raggedright\hyphenpenalty=10000}
  [\vspace{0.7ex}{\color{gold}\titlerule}]
\titlespacing*{\chapter}{0pt}{-10pt}{25pt}
\titleformat{\section}
  {\Large\sffamily\bfseries\color{navy}\raggedright\hyphenpenalty=10000}
  {\color{teal}\thesection}{0.7em}{}
\titleformat{\subsection}
  {\large\sffamily\bfseries\color{ink}}
  {\color{teal}\thesubsection}{0.7em}{}

\usepackage{fancyhdr}
\pagestyle{fancy}
\fancyhf{}
\fancyhead[LE]{\sffamily\small\color{slate}\thepage\quad\parbox[t]{0.90\textwidth}{\raggedright\nouppercase{\leftmark}}}
\fancyhead[RO]{\sffamily\small\color{slate}\parbox[t]{0.90\textwidth}{\raggedleft\nouppercase{\rightmark}}\quad\thepage}
\renewcommand{\headrulewidth}{0.45pt}
\renewcommand{\headrule}{\hbox to\headwidth{\color{teal}\leaders\hrule height \headrulewidth\hfill}}
\setlength{\headheight}{34pt}

\usepackage[hidelinks,bookmarksopen=true,pdfstartview=FitH]{hyperref}
\usepackage{bookmark}
\hypersetup{
  pdftitle={MAT0339 - Mathématiques, manuel de cours},
  pdfauthor={Xia Xiao},
  pdfsubject={Manuel structuré et illustré du cours MAT0339},
  pdfkeywords={MAT0339, mathématiques, algèbre, fonctions, probabilités, vecteurs, matrices, trigonométrie}
}

% ------------------------------------------------------------
% Raccourcis
% ------------------------------------------------------------
\newcommand{\R}{\mathbb{R}}
\newcommand{\N}{\mathbb{N}}
\newcommand{\Z}{\mathbb{Z}}
\newcommand{\Q}{\mathbb{Q}}
\newcommand{\C}{\mathbb{C}}
\newcommand{\vect}[1]{\overrightarrow{#1}}
\newcommand{\norm}[1]{\left\lVert #1\right\rVert}
\newcommand{\abs}[1]{\left\lvert #1\right\rvert}
\newcommand{\dom}{\operatorname{Dom}}
\newcommand{\im}{\operatorname{Im}}
\newcommand{\e}{\mathrm{e}}
\newcommand{\dd}{\,\mathrm{d}}

